\documentclass[11pt]{article}
\usepackage[utf8]{inputenc}
\usepackage{amsmath}
\usepackage{amssymb}
\usepackage{geometry}
\usepackage{enumitem}
\usepackage{xcolor}
\usepackage{hyperref}

\geometry{a4paper, margin=0.85in}
\hypersetup{colorlinks=true, linkcolor=blue!70!black, urlcolor=blue}

\usepackage{tcolorbox}
\tcbuselibrary{skins}
\newtcbox{\hl}{on line, arc=2pt, boxrule=0pt, colback=yellow!25!white, colframe=gray!50!black}

\title{\textbf{\color{red!70!black}EECS 498/598 Sample Exam Solutions}}
\author{Answer Key \& Grading Rubrics}
\date{\today}

\begin{document}
\maketitle

\section*{Question 1: Rasterization}

\paragraph{1a. Homogeneous coordinates allow translation to be represented as a matrix multiplication (same form as rotation and scaling).} Specifically, by adding the extra coordinate $w=1$, we can represent any affine transformation (Translation + Linear) using a single $3\times3$ (or $4\times4$ in 3D) matrix.

\paragraph{1b. Matrix multiplication is non-commutative ($A \cdot B \neq B \cdot A$).} For example, rotating a cup by $90^\circ$ around the origin then translating it to $(10,0)$ results in a different position than translating it to $(10,0)$ first and then rotating it by $90^\circ$ around the origin. In the second case, the translation vector itself gets rotated.

\paragraph{1c. Z-buffer stores the minimum depth ($z$-coordinate) seen so far for each pixel.} When a new fragment (potential pixel) arrives, its $z$ is compared to the stored value. If the new fragment is closer (smaller $z$, or larger depending on coordinate space convention), it overwrites the color buffer and updates the $z$-buffer. This solves visibility in an order-independent, hardware-efficient way ($O(1)$ per fragment).

\paragraph{1d. Bilinear filtering averages the 4 nearest texture pixels (texels), producing smooth transitions but still resulting in aliasing at a distance (Moiré patterns).} Anisotropic filtering takes into account the angle at which the texture is viewed. It addresses the fact that when viewing a floor stretching into the distance, the texture footprint on screen is a long narrow trapezoid, not a square. Anisotropic filtering samples more texels along the axis of elongation, preserving detail on surfaces viewed at steep angles.

\section*{Question 2: Ray Tracing}

\paragraph{2a. Rendering Equation:}
\[
L_o(\mathbf{x}, \omega_o) = L_e(\mathbf{x}, \omega_o) + \int_{\Omega} f_r(\mathbf{x}, \omega_i, \omega_o) L_i(\mathbf{x}, \omega_i) (\mathbf{n} \cdot \omega_i) \, d\omega_i
\]
The integral sums up all incoming light ($L_i$) reflected toward the viewer ($\omega_o$), weighted by the material (BRDF) and the angle (cosine term). It represents the infinite number of possible light directions $\omega_i$ hitting the point. Monte Carlo sampling is used because evaluating the integral analytically is usually impossible; we approximate it by averaging the result of random sampled directions.

\paragraph{2b. Recursive Ray Tracing spawns secondary rays at a hit point.} These include: (1) \textbf{Reflection rays} (direction determined by surface normal and incoming angle), (2) \textbf{Refraction rays} (bent direction based on material's index of refraction), and (3) \textbf{Shadow rays} (pointing toward light sources to check for occlusion). The final color is a weighted sum of these rays based on material properties.

\paragraph{2c. BVH avoids checking every triangle ($O(N) \rightarrow O(\log N)$).} It groups triangles into bounding boxes (nodes). When a ray is traced, we first check if it intersects the large Root Box. If not, we discard the entire subtree immediately (since the ray can't hit what's inside the box). If it does hit, we recurse down the children. This allows us to quickly skip empty space or distant objects.

\paragraph{2d. Radiance is "power per unit projected area per unit solid angle" ($W/m^2 \cdot sr$).} It characterizes the light traveling along a \textbf{specific direction}. Irradiance is "power per unit area" ($W/m^2$) arriving at a surface, integrating light coming from \textbf{all directions} in the hemisphere. The key difference is directionality: Radiance describes a beam/light ray; Irradiance describes the total accumulation of energy landing on a patch.

\section*{Question 3: Neural Fields}

\paragraph{3a. Advantages of Neural Fields:} (1) \hl{Resolution Independence}: They represent a continuous function, so you can query any point $(x,y,z)$ at infinite resolution (unlike Voxels which are blocky or Point Clouds which have gaps). (2) \hl{Memory Efficiency}: Millions or billions of data points are compressed into a few megabytes of neural network weights.

\paragraph{3b. Spectral Bias:} MLPs tend to learn low-frequency functions (smooth, blurry shapes) easily, but struggle to learn high-frequency details (sharp textures, edges). \hl{Positional Encoding} maps the input coordinates to a higher-dimensional space using sine and cosine functions of exponentially increasing frequencies. This provides the network with the high-frequency components it needs to reconstruct sharp details.

\paragraph{3c. Density $\sigma$ is a material property independent of sampling distance.} Opacity $\alpha$ depends on the step size (how far apart our samples are). By predicting $\sigma$, the formula remains consistent regardless of how many points the ray samples. They are related by:
\[
\alpha_i = 1 - \exp(-\sigma_i \cdot \delta_i)
\]
where $\delta_i$ is the distance between samples.

\paragraph{3d. Coarse-to-Fine Sampling:} The \hl{Coarse Network} uses uniform sampling to roughly estimate where density exists along the ray. It predicts a PDF (probability distribution) of where important features (high density/surfaces) are located. The \hl{Fine Network} then uses Importance Sampling to pick more points in those high-density regions (surfaces), resulting in sharper edges without increasing total computation.

\section*{Question 4: Geometry}

\paragraph{4a. Implicit Surfaces (SDF) vs Parametric (Meshes):}
\hl{Implicit} is defined by a function $f(x,y,z)=0$. It is \hl{better for checking inside/outside} because you simply plug in the point coordinates. If $f < 0$, it's inside. If $f > 0$, it's outside. Meshes require complex ray-casting or barycentric coordinate checks to determine containment.

\paragraph{4b. Bézier Curves Properties:}
The curve passes through the \hl{first and last control points} but does \hl{NOT pass through the middle control points}; it is merely tangent to them. The curve always remains strictly within the convex hull formed by all control points.

\paragraph{4c. Marching Cubes:}
It divides space into a grid of small cubes. It samples the implicit function at the \hl{8 corners} of each cube. If some corners are positive (outside) and some are negative (inside), the surface must pass through the cube. It calculates the intersection points on the edges (where function value is 0) and looks up a pre-defined table to generate the corresponding triangles connecting those points.

\section*{Question 5: Reconstruction}

\paragraph{5a. Camera Projection Equation:}
\[ \mathbf{x} = K \cdot [R \mid t] \cdot \mathbf{X} \]
\begin{itemize}
    \item \textbf{Intrinsics $K$}: Focal length ($f_x, f_y$), Principal point ($c_x, c_y$), skew. These are internal to the camera.
    \item \textbf{Extrinsics $[R \mid t]$}: Rotation matrix and translation vector. These describe the camera's position and orientation in the 3D world.
\end{itemize}

\paragraph{5b. Epipolar Constraint:}
Given a point $\mathbf{x}$ in Image 1, the corresponding point $\mathbf{x}'$ in Image 2 must lie on a specific line called the \textbf{Epipolar Line} ($\mathbf{l}' = F\mathbf{x}$). 
\hl{Simplification}: Instead of searching the entire 2D image (all $W \times H$ pixels) for a match, we only search along this 1D line, drastically reducing computation and ambiguity.

\paragraph{5c. MVS Limitations:}
MVS struggles with: (1) \hl{Textureless/Smooth surfaces} (like a white wall) because there are no features to match. (2) \hl{Specular/Reflective surfaces} (like a mirror or shiny metal) because the appearance changes depending on the view angle. (3) Occlusions.

\section*{Question 6: Generative Models}

\paragraph{6a. VAE vs GAN:}
\hl{Main Advantage of GANs}: They produce much sharper, higher-quality images because the adversarial loss forces the generator to fool the discriminator with realistic high-frequency details. VAEs often produce blurry images because they use a pixel-wise MSE (L2) loss, which tends to average over possible outcomes.

\paragraph{6b. Diffusion Process:}
\hl{Forward}: We take a real image and gradually add Gaussian noise to it over $T$ steps until it becomes pure random noise (destroying all structure).
\hl{Reverse}: The neural network is trained to predict the added noise at each step. By iteratively removing this predicted noise, the network can reconstruct a clean image starting from pure random noise.

\paragraph{6c. Classifier-Free Guidance:}
It allows us to control generation strength without training a separate classifier. At inference time, we combine the output of the conditional prediction ($\epsilon_c$, e.g., "a cat") and the unconditional prediction ($\epsilon_u$, no condition).
Formula:
\[ \hat{\epsilon} = \epsilon_u + w \cdot (\epsilon_c - \epsilon_u) \]
Higher $w$ means stronger adherence to the prompt.

\paragraph{6d. Diffusion vs Flow Models:}
\hl{Diffusion} relies on a stochastic differential equation (SDE); the reverse process involves adding predicted noise in many small stochastic steps. 
\hl{Flow Models (Flow Matching)} learn a deterministic velocity field (ODE). They define a direct path (ODE) from noise to data. This allows for more direct, efficient integration with fewer steps compared to the hundreds or thousands required by Diffusion models.

\end{document}